\documentclass[border=5pt]{standalone}
\usepackage{circuitikz}
\begin{document}
\begin{circuitikz}[american]
  % izquierda: V_s en serie con R
  \draw (0,0) to[V=$V_s$] (0,3) to[R=$R$] (2.5,3) -- (3,3) node[right]{$a$};
  \draw (0,0) -- (3,0) node[right]{$b$};
  \node at (1.5,-1){tensión real};
  % flecha equivalencia
  \node at (4.6,1.5){\Large $\equiv$};
  % derecha: I_s en paralelo con R
  \draw (6,0) to[I=$I_s$] (6,3) -- (7.5,3);
  \draw (7.5,3) to[R=$R$] (7.5,0);
  \draw (6,0) -- (7.5,0);
  \draw (7.5,3) -- (9,3) node[right]{$a$};
  \draw (7.5,0) -- (9,0) node[right]{$b$};
  \node at (7.5,-1){corriente real};
\end{circuitikz}
\end{document}
